\ifdefined\HMTTRITFormal
\chapter{Estado local ternario y orientación del transporte}
\label{chap:trit}
\index{TRIT}

El TRIT es la coordenada ternaria que conserva la orientación local de una
transición. Su definición precede a cualquier interpretación temporal: los
tres valores distinguen tipos algebraicos de transporte y sólo adquieren una
lectura de pasado, umbral o futuro cuando se fija un parámetro ordenado.

\section{Firma ternaria orientada y espacio de estados}

El espacio de firmas visibles es
\[
\TRIT=\{-1,0,+1\}.
\]
Sus elementos son símbolos orientados, no resultados de una operación en
\(\ZZ\). El mapa balanceado
\[
\FF_3\longrightarrow\TRIT,
\qquad0\mapsto0,\quad1\mapsto+1,\quad2\mapsto-1,
\]
permite pasar del residuo ternario a la firma orientada.

\begin{definition}[Estado completo del TRIT y fibra de acarreo]
\label{def:trit-levantado-carry-rev8}
En un paso local cuya amplitud requiere a lo sumo un acarreo, la
descomposición ternaria balanceada se escribe
\[
 \widehat\tau=(r,c),\qquad
 a+b=r+3c,
 \qquad r,c\in\{-1,0,+1\}.
\]
La primera coordenada \(r\) es el residuo visible y la segunda \(c\) conserva
el cociente que hizo posible esa lectura.  En particular, la fibra sobre el
cero visible contiene tres estados distintos:
\[
 0^+=(0,+1),\qquad 0^0=(0,0),\qquad 0^-=(0,-1),
 \qquad
 \operatorname{vis}(0^+)=\operatorname{vis}(0^0)
 =\operatorname{vis}(0^-)=0.
\]
\end{definition}

Esta separación local entre residuo y cociente es la que, en la elevación
Bockstein--Hensel posterior, adopta la forma matricial
\[
 x_k A=u_k+3x_{k+1}.
\]
Allí \(u_k\) publica el residuo del nivel \(k\), mientras \(x_{k+1}\)
transporta la memoria necesaria para recuperar el levantamiento \(x_k\).  Por tanto,
un cero visible no determina por sí solo el estado ni su futuro admisible.
Esta afirmación es algebraica: no convierte \(0^\pm\) en un infinitesimal,
una energía, una torsión física ni un intervalo temporal.

Un trit visible no determina el conjunto de continuaciones admisibles. El
espacio completo de estados de transporte contiene los campos
\[
x=(t,\mu,o,h,c,\sigma,\lambda,g).
\]
Aquí \(t\in\TRIT\) es la firma visible;
\(\mu\in\mathsf M:=\{\Sigma,\Pi,\varnothing\}\), la etiqueta del observable
aditivo, multiplicativo o neutro; \(o\), la orientación; \(h\), la coordenada de extensión;
\(c\), el acarreo acumulado; \(\sigma\), la firma de frontera; \(\lambda\), el
registro de trayectoria; y
\(g\in\ZZ/9\ZZ\), la fase. La proyección \(x\mapsto t\) elimina todas las
demás componentes.

\begin{remark}[Distinción entre símbolo y estado]
Dos estados pueden mostrar el mismo trit y poseer continuaciones distintas.
Una palabra ternaria tampoco determina por sí sola la etiqueta de observable,
la orientación, el acarreo o la
monodromía. Esta distinción explica por qué cinco regiones de \(\pi\) pueden
compartir una palabra \(w_\pi\) sin convertirse en una sola región.
\end{remark}

\fi

\ifdefined\HMTTRITDevelopments
\chapter{Álgebras cuadráticas y regímenes de transporte}
\label{chap:trit-algebras}

\section{\(3\) álgebras cuadráticas asociadas}

Para \(\tau\in\TRIT\), sea \(\bar\tau\in\{-1,0,+1\}\subset\RR\) su
imagen por la identificación balanceada anterior. Consideramos el álgebra real
unitaria
\[
\mathbb A_\tau=\RR[X]/(X^2+\bar\tau),
\]
en la que \(J\) denota la clase de \(X\) e \(I\) la unidad. Esta convención
separa el símbolo TRIT de su valor escalar.
Según el signo se obtienen tres tipos:
\[
\begin{array}{ccl}
\tau=+1&:&J^2=-I,\quad\text{tipo elíptico o complejo},\\
\tau=0&:&J^2=0,\quad\text{tipo parabólico o nilpotente},\\
\tau=-1&:&J^2=+I,\quad\text{tipo hiperbólico o descompuesto}.
\end{array}
\]
Las exponenciales exhiben la diferencia:
\[
e^{tJ}=
\begin{cases}
\cos t+J\sin t,&J^2=-I,\\
1+tJ,&J^2=0,\\
\cosh t+J\sinh t,&J^2=+I.
\end{cases}
\]
Las tres álgebras son, respectivamente, la extensión compleja, el álgebra de
números duales y el álgebra real descompuesta. Sus subgrupos exponenciales
producen evoluciones elípticas, parabólicas e hiperbólicas sin introducir una
geometría continua adicional en la definición del TRIT\@.
Las realizaciones reales concretas seleccionadas por el Coxeter de APP, el
generador nilpotente y la incidencia de hojas TPK se construyen, sin
duplicar esta clasificación abstracta, en
\cref{sec:tres-generadores-concretos-trit}.

\begin{figure}[H]
\centering
\resizebox{\textwidth}{!}{\begin{tikzpicture}[font=\small,box/.style={draw,rounded corners=2mm,minimum width=4.2cm,minimum height=4.4cm,align=center,inner sep=3mm}]
  \node[box,draw=hmtviolet,fill=hmtlightviolet] (ell) at (0,0) {};
  \node[font=\bfseries,align=center,text width=3.6cm,hmtviolet]
    at (0,1.55) {Elíptica /\\compleja};
  \draw[hmtviolet,thick,->] (0,-.25) arc (-90:255:.62);
  \node[fill=hmtlightviolet,inner sep=1.5pt] at (0,.32) {\(J^2=-I\)};
  \node[align=center] at (0,-1.18) {subgrupo compacto\\retorno elíptico};

  \node[box,draw=hmtgray,fill=hmtpaper] (par) at (5,0) {};
  \node[font=\bfseries,align=center,text width=3.6cm,hmtgray]
    at (5,1.55) {Parabólica /\\nilpotente};
  \draw[hmtgray,thick,->] (4.1,.3) .. controls (4.6,-.2) and (5.4,-.2) .. (5.9,.3);
  \draw[hmtgray,dashed] (4.05,.3)--(5.95,.3);
  \node at (5,.72) {\(J^2=0\)};
  \node[align=center] at (5,-1.18) {flujo unipotente\\estado umbral};

  \node[box,draw=hmtgold,fill=hmtlightgold] (hyp) at (10,0) {};
  \node[font=\bfseries,align=center,text width=3.6cm,hmtgold]
    at (10,1.55) {Hiperbólica /\\desdoblada};
  \draw[hmtgold,thick,->] (9.35,-.55) .. controls (9.62,-.10) and (9.62,.35) .. (9.78,.60);
  \draw[hmtgold,thick,->] (10.65,-.55) .. controls (10.38,-.10) and (10.38,.35) .. (10.22,.60);
  \node[fill=hmtlightgold,inner sep=1.5pt] at (10,.05) {\(J^2=+I\)};
  \node[align=center] at (10,-1.18) {dos subespacios propios\\evolución hiperbólica};

  \node[font=\footnotesize,align=center,hmtgray] at (5,-2.75) {La interpretación temporal se obtiene al fijar una parametrización ordenada de los tres tipos algebraicos.};
\end{tikzpicture}
}
\caption{Tres tipos algebraicos asociados al estado ternario orientado. La
primera línea recoge las identidades algebraicas; la segunda, su interpretación
respecto de una parametrización ordenada.}
\label{fig:trit-atlas}
\end{figure}

\section{Compatibilidad y trayectorias parametrizadas}

La expresión \emph{All Possible Paths} invita a imaginar caminos recorridos en
el tiempo. HMT sitúa el objeto un paso antes. Sea \(\Omega_{\rm HMT}\) un
conjunto de estados locales provisto de una relación de transición admisible
\(\leadsto\), una firma TRIT y datos de transporte que se tiparán en el
apartado siguiente. Si
\(\mathcal T\) es un conjunto discretamente ordenado, una trayectoria
parametrizada es una aplicación
\[
\gamma:\mathcal T\longrightarrow\Omega_{\rm HMT}
\]
que satisface \(\gamma(t)\leadsto\gamma(t^+)\) para cada sucesor \(t^+\).
El orden de \(\mathcal T\) convierte la compatibilidad en una sucesión
parametrizada. Antes de elegir un orden sólo hay una familia compatible
indexada por el grafo o poset correspondiente; no se la denomina sección salvo
que se haya declarado una proyección \(p:E\to\mathcal T\) con
\(p\circ\gamma=\operatorname{id}_{\mathcal T}\).
En consecuencia, una suma de caminos finita puede leerse como suma sobre
temporalizaciones de una estructura que, en su nivel primario, no presupone
un tiempo exterior.

\section{Parametrización ordenada de la compatibilidad}

Una parametrización ordenada permite interpretar retorno, umbral y apertura
como tres posiciones relativas en una evolución. La convención utilizada en la
\cref{fig:trit-atlas} es
\[
\begin{array}{ccl}
J^2=-I&\rightsquigarrow&\text{retorno y orientación inversa},\\
J^2=0&\rightsquigarrow&\text{estado umbral},\\
J^2=+I&\rightsquigarrow&\text{descomposición en dos direcciones propias}.
\end{array}
\]
En la parametrización temporal adoptada, \(+1\) representa
retorno, memoria y pasado; \(0\), umbral y presente; y \(-1\), apertura
bidireccional y futuro. Esta lectura ordenada no altera los tres tipos
cuadráticos precedentes.
La involución exacta, tipada para cada \(n\geq1\),
\[
C_{\rm rev}:\TRIT^n\longrightarrow\TRIT^n,
\qquad
C_{\rm rev}(\tau_1,\ldots,\tau_n)=(-\tau_n,\ldots,-\tau_1)
\]
intercambia \(+1\) y \(-1\), preserva \(0\) y satisface
\(C_{\rm rev}^2=\operatorname{id}_{\TRIT^n}\).
Así, la orientación temporal no se introduce como un cuarto estado: resulta
de la parametrización de la firma junto con una involución que conserva el
umbral.

\section{Descomposición espectral del tipo hiperbólico}

En el tipo hiperbólico, los proyectores
\[
P_\pm=\frac{I\pm J}{2}
\]
separan dos direcciones propias. La involución que cambia la orientación las
intercambia. Algebraicamente, \(P_+P_-=0\), \(P_++P_-=I\) y
\(JP_\pm=\pm P_\pm\). Por tanto el tipo hiperbólico contiene dos subespacios
propios complementarios y una operación exacta que los permuta.

\begin{remark}[Estructura completa del estado TRIT]
TRIT no es un sustituto de los signos \(-,0,+\). Es el observable mínimo de
un estado con etiqueta de observable, orientación y registro de trayectoria.
Sus tres álgebras asociadas hacen posible
que un mismo soporte admita cierres elípticos, umbrales parabólicos y aperturas
hiperbólicas antes de fijar una evaluación temporal.
\end{remark}
\fi
